本容量是 $n=100$, 我们应当预期有 5 个观测值的 $d_{j}^{2}$ 值超过卡方分布的上第 5 个百分位数. 一个较为极端的百分位数应用于确定不拟合其余数据的图形的观测值.

\hyperref[tab:4.3]{表 4.3} 中给出的关于木板的数据已经进行过一些清理. 来自同一研究的类似数据集还包含对 $x_{5}=$张力强度的数据. 下表的各行列出的是从 112 个观测值向量中选出的 9 个向量, 并给出了其标准化值.

\begin{mvtable}
\centering
\begin{tabular}{ccccc|ccccc}
\toprule
$x_{1}$ & $x_{2}$ & $x_{3}$ & $x_{4}$ & $x_{5}$ & $z_{1}$ & $z_{2}$ & $z_{3}$ & $z_{4}$ & $z_{5}$ \\
\midrule
$\vdots$ & $\vdots$ & $\vdots$ & $\vdots$ & $\vdots$ & $\vdots$ & $\vdots$ & $\vdots$ & $\vdots$ & $\vdots$ \\
1\,631 & 1\,528 & 1\,452 & 1\,559 & 1\,602 & 0.06 & -0.15 & 0.05 & 0.28 & -0.12 \\
1\,770 & 1\,677 & 1\,707 & 1\,738 & 1\,785 & 0.64 & 0.43 & 1.07 & 0.94 & 0.60 \\
1\,376 & 1\,190 & 723 & 1\,285 & 2\,791 & -1.01 & -1.47 & -2.87 & -0.73 & \tikz[baseline=(a.base)]\node[draw,circle,inner sep=1pt](a){4.57}; \\
1\,705 & 1\,577 & 1\,332 & 1\,703 & 1\,664 & 0.37 & 0.04 & -0.43 & 0.81 & 0.13 \\
1\,643 & 1\,535 & 1\,510 & 1\,494 & 1\,582 & 0.11 & -0.12 & 0.28 & 0.04 & -0.20 \\
1\,567 & 1\,510 & 1\,301 & 1\,405 & 1\,553 & -0.21 & -0.22 & -0.56 & -0.28 & -0.31 \\
1\,528 & 1\,591 & 1\,714 & 1\,685 & 1\,698 & -0.38 & 0.10 & 1.10 & 0.75 & 0.26 \\
1\,803 & 1\,826 & 1\,748 & 2\,746 & 1\,764 & 0.78 & 1.01 & 1.23 & \tikz[baseline=(b.base)]\node[draw,circle,inner sep=1pt](b){4.65}; & 0.52 \\
1\,587 & 1\,554 & 1\,352 & 1\,554 & 1\,551 & -0.13 & -0.05 & -0.35 & 0.26 & -0.32 \\
$\vdots$ & $\vdots$ & $\vdots$ & $\vdots$ & $\vdots$ & $\vdots$ & $\vdots$ & $\vdots$ & $\vdots$ & $\vdots$ \\
\bottomrule
\end{tabular}
\end{mvtable}

标准化值是根据从所有 112 个观测值计算出的样本均值和方差得到的. 表中有两个极端标准化值. 这两个值都太大, 超过 4.5. 在研究时, 研究人员是在记录本中用手记录测量值, 然后进行计算, 得出表中的值. 当他们根据分析结果核对这两个值的记录时, 发现了错误所在. 值 $x_{5}=2791$ 被改为 1241, $x_{4}=2746$ 被改为 1670. 通过找到标准化值的大约首位数字的位置, 就可以迅速地发现在个别变量上的错误读数.

下面的例题返回到例 4.14 中讨论过的木板数据.

\noindent\textbf{例 4.15}\quad \textbf{(在木板数据中搜寻离群值)}

\hyperref[tab:4.4]{表 4.4} 包含了\hyperref[tab:4.3]{表 4.3} 的数据以及标准化观测值. 这些数据由硬度的四种不同的测量值 $x_{1},x_{2},x_{3}$ 和 $x_{4}$ 组成, 每种都有 $n=30$ 块木板. 第一个测量值是向下传送一个冲击波给木板, 第二个测量值是当木板振动时确定的, 最后两个测量值是从静态的试验得到的. 标准化的测量值是
\[
z_{jk}=\frac{x_{jk}-\bar{x}_{k}}{\sqrt{s_{kk}}},\quad k=1,2,3,4;\quad j=1,2,\cdots,30
\]
距离的平方是 $d_{j}^{2}=(\mathbf{x}_{j}-\bar{\mathbf{x}})'\mathbf{S}^{-1}(\mathbf{x}_{j}-\bar{\mathbf{x}})$.

\hyperref[tab:4.4]{表 4.4} 中最后一栏显示出样本 16 是一个多元离群值, 因为 $\chi_{4}^{2}(0.005)=14.86$; 所有的单个测量值都处于它们各自的单变量散布范围之内. 样本 9 也有一个大的 $d^{2}$ 值.

这两个平方距离较大的样本(9 和 16), 十分明显与\hyperref[fig:4.9]{图 4.9} 中的其余部分不同. 一旦这两点被除去, 剩余部分的图形符合期望的直线关系. 对于木板硬度测量的散布图, 在\hyperref[fig:4.11]{图 4.11} 中给出. 图中的实心圆点对应于样本 9 和 16. 虽然样本 16 的点在所有图中都与众不同, 样本 9 的点在 $x_{3}$ 与 $x_{1}$ 的散布图中是“隐藏”的, 并且在 $x_{1}$ 与 $x_{2}$ 的散布图中几乎是隐藏的. 然而, 当所有四个变量都考虑时, 样本 9 很清楚地确定为一个多元离群值.
